\documentclass[UTF8,zihao=-4,fontset=fandol]{ctexart}
\usepackage[a4paper,margin=22mm,headheight=16pt]{geometry}
\usepackage{amsmath,amssymb,booktabs,longtable,array,tabularx}
\usepackage{xcolor,fancyhdr,enumitem,fvextra,graphicx}
\usepackage[unicode,colorlinks=true,linkcolor=teal!65!black,urlcolor=teal!65!black]{hyperref}
\definecolor{Ink}{HTML}{17324D}
\definecolor{Accent}{HTML}{007F83}
\definecolor{CodeBG}{HTML}{F2F5F7}
\setlength{\parindent}{0pt}
\setlength{\parskip}{5pt}
\setlength{\emergencystretch}{3em}
\setlist{itemsep=2pt,topsep=3pt,leftmargin=1.6em}
\ctexset{section={format=\Large\bfseries\color{Ink}},subsection={format=\large\bfseries\color{Accent}}}
\pagestyle{fancy}
\fancyhf{}
\fancyhead[L]{\small ExperimentAgent / Lecture 01}
\fancyhead[R]{\small 2026-10-21}
\fancyfoot[C]{\small\thepage}
\DefineVerbatimEnvironment{Code}{Verbatim}{fontsize=\small,breaklines=true,breakanywhere=true,frame=single,framesep=6pt,rulecolor=\color{Accent!30},formatcom=\color{Ink}}
\newcommand{\code}[1]{{\small\ttfamily\detokenize{#1}}}
\hypersetup{pdftitle={Agent 学习讲义 01：从文件工具到双工具真实循环},pdfauthor={Agent Study}}
\begin{document}
\hypersetup{pageanchor=false}
\begin{titlepage}
\vspace*{14mm}
{\large\color{Accent}实验驱动学习 / 第一阶段复盘\par}
\vspace{10mm}
{\Huge\bfseries\color{Ink}Agent 学习讲义 01\par}
\vspace{6mm}
{\LARGE 从文件工具到双工具真实循环\par}
\vspace{12mm}
{\large 概念、代码、实验、错题与复习路线\par}
\vspace{14mm}
\hrule
\vspace{6mm}
学习日期：2026 年 10 月 21 日\\
项目：ExperimentAgent（代码实验助手）\\
学习环境：Windows / Anaconda / Python 3.11 / 本地 Ollama\\
本讲范围：16 次小课、14 个运行实验、1 次独立功能修改\\
当前关卡：G1 前的工程基础；已跑通最小模型与工具闭环
\vfill
\small
本文是本项目的阶段性学习讲义，不是 CS2680 官方讲义。\\
实际运行、教学模拟、独立修改与后续规划分别标记。\\
原始源码、笔记和报告通过附录索引保留，复盘数字以对应运行记录为准。
\end{titlepage}
\hypersetup{pageanchor=true}
\pagenumbering{roman}
\tableofcontents
\clearpage
\pagenumbering{arabic}

学习日期：2026-10-21。项目：ExperimentAgent，代码实验助手。适用环境：Windows、Anaconda、Python 3.11、本地 Ollama。

本讲覆盖本项目第 1–16 次小课、l01–l14 实验及一次独立功能修改。它是阶段复盘，不是 CS2680 官方讲义；课程讲义和论文尚未逐项阅读验收。当前主要覆盖 16 周路线前 1–3 周的工程内容，尚未整体验收 G1。

公开版说明：用户目录已用 USER 占位符脱敏，附件内部路径与无关环境名称已移除；实验结论、错误类别和耗时保留。快照清单的字节数和 SHA-256 对应公开脱敏版本。

\section{先建立一张完整的知识地图}

\subsection{今天已经搭出了什么}

现在的程序能接收任务，将工具说明交给本地模型，由模型提出工具名称与参数；Python 检查、执行并记录真实结果；模型根据观察继续行动或回答。程序具备步数上限、参数修正预算和部分文件异常的结构化处理。

当前工具只有两种：\nolinkurl{read_text_file} 读取 UTF-8 正文；\nolinkurl{list_directory} 查看目录的一层条目。还没有运行训练、编辑代码、持久化记忆、语义检索或实验指标验证。

\begin{Code}
用户目标
  → 模型读取任务、历史和工具说明
  → assistant 提出 tool_calls
  → Python 校验名称、参数与路径范围
  → 工具实际访问文件系统
  → 记录 role=tool 的真实结果
  → 策略判断继续、修正或停止
  → 再次请求模型，或者由运行时终止
\end{Code}

模型提出动作，工具接触环境，运行时管理消息与执行边界。模型说“已经读取”不能替代真实读取结果；运行时返回 completed 也不自动证明任务内容正确。

\subsection{Chat、workflow 与 Agent}

{\small\setlength{\tabcolsep}{4pt}
\renewcommand{\arraystretch}{1.25}
\begin{longtable}{>{\raggedright\arraybackslash}p{0.23\linewidth}>{\raggedright\arraybackslash}p{0.33\linewidth}>{\raggedright\arraybackslash}p{0.33\linewidth}}
\toprule
\textbf{形式} & \textbf{下一步由谁决定} & \textbf{本项目中的例子} \\
\midrule\endfirsthead
\toprule
\textbf{形式} & \textbf{下一步由谁决定} & \textbf{本项目中的例子} \\
\midrule\endhead
\bottomrule\endfoot
单轮聊天 & 一次模型生成 & l08 用一句话解释工具调用 \\
固定工作流 & 程序预先安排步骤 & 先运行 l09，再运行 l10 回传结果 \\
模型驱动循环 & 模型结合历史与观察提出动作，程序负责执行 & l11 自动完成读取；l14 在两种工具间选择 \\
\end{longtable}}

一个任务即使使用了模型，也可能仍是固定流程。l14 的组合任务本身要求先看目录再读取；这验证了这一小范围内的工具选择与观察使用，不代表通用规划能力已经通过评测。

\subsection{三类证据要分开}

{\small\setlength{\tabcolsep}{4pt}
\renewcommand{\arraystretch}{1.25}
\begin{longtable}{>{\raggedright\arraybackslash}p{0.23\linewidth}>{\raggedright\arraybackslash}p{0.33\linewidth}>{\raggedright\arraybackslash}p{0.33\linewidth}}
\toprule
\textbf{证据} & \textbf{能支持什么} & \textbf{不能替代什么} \\
\midrule\endfirsthead
\toprule
\textbf{证据} & \textbf{能支持什么} & \textbf{不能替代什么} \\
\midrule\endhead
\bottomrule\endfoot
预设 mock 请求 & 循环在指定失败输入下行为可控 & 真实模型的纠错或推理能力 \\
一次真实模型运行 & 该次任务中的实际调用与回答 & 任意任务上的可靠性结论 \\
学习者独立修改并通过验收 & 这一具体分支的实现能力 & 从空白编写整个 Agent 的能力 \\
\end{longtable}}

\section{环境与复现方式}

\subsection{使用专用 Anaconda 环境}

在 Anaconda Prompt 中运行：

\begin{Code}
cd /d E:\study\agent
conda activate agent-study
python -c "import sys; print(sys.executable)"
\end{Code}

已确认解释器为 E:/Anaconda3/envs/agent-study/python.exe，Python 3.11.17。Agent 学习代码目前使用 Python 标准库，没有要求安装完整 Agent 框架。

\begin{Code}
python -X utf8 -m labs.l01_read_file
python -X utf8 -m unittest discover -s tests -v
\end{Code}

\begin{itemize}
\item \code{-X utf8}：启用 Python UTF-8 模式，便于 Windows 终端和文本处理。
\item \code{-m labs.xxx}：从项目根目录按模块运行，便于使用项目内导入。
\item \code{discover -s tests}：发现 tests 目录中的 unittest 测试。
\end{itemize}

\code{environment.yml} 记录了 python=3.11 与 pip。原有环境不需要反复重建。后续安装新依赖时应单独记录版本，不把“已安装”写成“已验证可用”。

\subsection{本机模型信息}

{\small\setlength{\tabcolsep}{4pt}
\renewcommand{\arraystretch}{1.25}
\begin{longtable}{>{\raggedright\arraybackslash}p{0.28\linewidth}>{\raggedright\arraybackslash}p{0.63\linewidth}}
\toprule
\textbf{项目} & \textbf{已观察到的配置} \\
\midrule\endfirsthead
\toprule
\textbf{项目} & \textbf{已观察到的配置} \\
\midrule\endhead
\bottomrule\endfoot
本地 API & http://127.0.0.1:11434/api/chat \\
Ollama 版本 & 0.40.6，来自本机元数据记录 \\
主实验模型 & qwen3.5:9b-\nolinkurl{q4_K_M} \\
另一个已安装模型 & qwen3.5:2b \\
GPU & RTX 5080，总显存 16303 MiB，先前只读检查记录 \\
当前运行接口 & urllib.request 的 JSON HTTP 请求 \\
\end{longtable}}

这些是本次机器与运行的记录，不是推荐的最新软件版本。权重安装和元数据查询不等于模型推理成功；聊天成功也不等于工具调用已成功。

\section{从普通文件读取函数开始}

\subsection{先掌握输入、动作与返回值}

\begin{Code}
path = Path(project_root) / relative_path
content = path.read_text(encoding="utf-8")
return {"status": "success", "content": content}
\end{Code}

参数告诉工具读取哪里，\nolinkurl{read_text} 才真正访问文件。返回的字典不是模型回答，而是 Python 执行得到的工具观察。

已有文件实际内容为：

\begin{Code}
你好，Agent！
这是一次真实的文件读取。
\end{Code}

正文的最后一个换行也属于内容；精确测试会检查它。\nolinkurl{read_text} 的成功并不等于后续模型已经看见结果。

\subsection{路径为什么容易出错}

最初 train.csv 场景是纸面练习，不是已经运行过的训练。若进程 cwd 是 E:/study/agent，普通相对路径 data/train.csv 对应 E:/study/agent/data/train.csv。

如果数据实际在 labs/demo/data/train.csv，可在该 cwd 下修正相对路径，或者显式以正确根目录为锚点。我们的工具采用第二种思路：根目录由运行程序持有，模型只提供 \nolinkurl{relative_path}。

FileNotFoundError 只表明在实际解析位置没有找到目标。还要区分数据未准备、路径写错、cwd 不符等原因；不能因为缺文件就擅自生成训练数据。

\subsection{先检查授权范围，再访问环境}

\begin{Code}
root = Path(project_root).resolve()
path = (root / relative_path).resolve()
if not path.is_relative_to(root):
    return {"status": "failed", "error_type": "path_outside_project"}
\end{Code}

resolve 将路径转换为规范化路径，并解析可解析的链接；\nolinkurl{is_relative_to} 判断它是否处于项目根目录下。不能用字符串前缀判断，例如 project-copy 与 project 可能共享文本前缀，却不是同一目录范围。

{\small\setlength{\tabcolsep}{4pt}
\renewcommand{\arraystretch}{1.25}
\begin{longtable}{>{\raggedright\arraybackslash}p{0.28\linewidth}>{\raggedright\arraybackslash}p{0.63\linewidth}}
\toprule
\textbf{输入或环境} & \textbf{应发生的事} \\
\midrule\endfirsthead
\toprule
\textbf{输入或环境} & \textbf{应发生的事} \\
\midrule\endhead
\bottomrule\endfoot
项目内存在的文件 & 允许访问，成功读取 \\
项目内不存在的文件 & 范围允许；实际读取返回 \nolinkurl{file_not_found} \\
../outside.txt & 访问前拒绝，\nolinkurl{path_outside_project} \\
项目外绝对路径，即使文件不存在 & 先拒绝越界，不先报缺文件 \\
\end{longtable}}

范围许可与目标存在是两个不同问题。当前检查不是完整操作系统沙箱，也没有解决所有并发路径变化与资源隔离问题。

\section{注册表：名称如何变成函数调用}

\subsection{字符串、函数对象和函数返回值}

\begin{Code}
tools = {"read_text_file": read_project_file}
tool = tools["read_text_file"]
result = tool(relative_path="labs/fixtures/hello.txt")
\end{Code}

\begin{itemize}
\item \nolinkurl{"read_text_file"} 是字典键，也就是公开工具名称。
\item \nolinkurl{read_project_file} 没有括号，表示函数对象。
\item \nolinkurl{read_project_file(...)} 带括号，表示现在执行函数，并取得返回值。
\end{itemize}

学习者已经正确复述：先取出函数，下一行才开始执行读取。注册时若写成函数调用，通常存进去的是那一次调用的返回值，不再是稍后可调用的函数。

\subsection{包装函数固定可信根目录}

\begin{Code}
root = str(Path(project_root).resolve())

def read_project_file(relative_path):
    return read_text_file(root, relative_path)
\end{Code}

包装函数记住构造注册表时的 root。不同注册表可以绑定不同项目根目录，互不混用；模型不能额外提供 \nolinkurl{project_root} 来改变授权范围。

\subsection{分发与参数展开}

\begin{Code}
tool = tools.get(name)
result = tool(**arguments)
\end{Code}

当 arguments 为 \code{{"relative_path": "hello.txt"}}，\code{**arguments} 相当于关键字调用 \nolinkurl{tool(relative_path="hello.txt")}。未知名称返回 \nolinkurl{unknown_tool}；不存在的删除工具不会因模型给出名称而自动获得实现。

\section{参数校验与 JSON Schema}

\subsection{参数校验在实际执行之前}

两个路径工具现在复用 \nolinkurl{validate_relative_path_arguments}，依次检查：

\begin{enumerate}
\item arguments 必须是字典。
\item 必须包含 \nolinkurl{relative_path}。
\item 不能带额外字段，包括 \nolinkurl{project_root}。
\item 值必须是字符串。
\item 字符串不能是空串，也不能全部为空白。
\end{enumerate}

\begin{Code}
value = arguments["relative_path"]
if not isinstance(value, str) or not value.strip():
    return "relative_path must be a non-empty string"
\end{Code}

这里利用 or 的短路求值：不是字符串时，右边不会调用 strip。校验器返回错误说明，分发器将其包装为 \nolinkurl{invalid_arguments}；参数不合法时，工具不执行。

\subsection{strip 的准确含义}

strip 只移除两端空白，返回一个新字符串，不会删除中间空格，也不会原地改写输入。本项目只用 strip 的结果判断是否全空白，实际传给工具的路径保持原样。

{\small\setlength{\tabcolsep}{4pt}
\renewcommand{\arraystretch}{1.25}
\begin{longtable}{>{\raggedright\arraybackslash}p{0.23\linewidth}>{\raggedright\arraybackslash}p{0.33\linewidth}>{\raggedright\arraybackslash}p{0.33\linewidth}}
\toprule
\textbf{字符串} & \textbf{校验是否通过} & \textbf{说明} \\
\midrule\endfirsthead
\toprule
\textbf{字符串} & \textbf{校验是否通过} & \textbf{说明} \\
\midrule\endhead
\bottomrule\endfoot
空字符串 & 否 & 没有内容 \\
三个空格 & 否 & strip 后为空 \\
\code{   a   } & 是 & 包含非空白字符；实际路径仍带两端空格 \\
\code{my notes.txt} & 是 & 中间空格是文件名的一部分 \\
\end{longtable}}

\subsection{把参数规则描述给模型}

\begin{Code}
{
  "type": "object",
  "properties": {
    "relative_path": {
      "type": "string",
      "minLength": 1,
      "pattern": "\\S"
    }
  },
  "required": ["relative_path"],
  "additionalProperties": false
}
\end{Code}

\begin{itemize}
\item minLength: 1 只检查长度；空格也计算长度。
\item pattern 的正则 \code{\S} 表示至少出现一个非空白字符；在 JSON 源码中反斜杠需要转义。
\item required 指定必填键；additionalProperties: false 描述不允许额外键。
\end{itemize}

\code{   a   } 满足这两个字符串条件。Schema 不会替我们执行 strip，也不会证明路径存在、位于项目内或具备访问权限。

Schema 是模型看到的输入契约；Python 校验负责真正执行约束。json.loads 只是解析 JSON，本项目尚未安装通用 JSON Schema 验证器。

\section{消息协议：模型怎样提出工具请求}

\subsection{四种角色}

{\small\setlength{\tabcolsep}{4pt}
\renewcommand{\arraystretch}{1.25}
\begin{longtable}{>{\raggedright\arraybackslash}p{0.28\linewidth}>{\raggedright\arraybackslash}p{0.63\linewidth}}
\toprule
\textbf{role} & \textbf{本项目中的用途} \\
\midrule\endfirsthead
\toprule
\textbf{role} & \textbf{本项目中的用途} \\
\midrule\endhead
\bottomrule\endfoot
system & 设置助手身份及读取必须来自工具等规则 \\
user & 用户任务或用户提供的资料，不只限于问句 \\
assistant & 模型给出的文本或工具调用请求 \\
tool & 实际工具执行结果，包括成功和失败 \\
\end{longtable}}

角色说明来源与用途；不能仅根据内容是不是问题来决定 role。工具失败时依然使用 role=tool。

\subsection{工具说明不是 Python 函数本体}

Ollama 请求中的 tools 使用以下外层包装：

\begin{Code}
{
  "type": "function",
  "function": {
    "name": "read_text_file",
    "description": "读取项目内指定 UTF-8 文件的文本内容。",
    "parameters": {"type": "object"}
  }
}
\end{Code}

这里 parameters 是缩略示意，实际发送完整 Schema。工具说明不包含可执行 Python 函数，也没有自动附带文件正文。模型返回名称与参数，由本地运行时查找函数。

\subsection{真实模型返回的请求}

\begin{Code}
{
  "role": "assistant",
  "content": "",
  "tool_calls": [{
    "id": "call_example",
    "function": {
      "name": "read_text_file",
      "arguments": {"relative_path": "labs/fixtures/hello.txt"}
    }
  }]
}
\end{Code}

示例 id 仅为说明；实际 id 在报告中。content 为空不一定失败，模型可能通过 \nolinkurl{tool_calls} 提出操作。done=true、\nolinkurl{done_reason}=stop 只表示本轮生成结束，不表示整个用户任务完成。

\nolinkurl{tool_calls} 是列表，一轮可能有多个请求。本程序会处理完整批次，不只取第一项；执行方式仍是顺序执行，没有并发工具执行。

\section{读取成功以后，怎样真正回传}

\subsection{保留完整 assistant 请求}

\begin{Code}
messages.append(assistant_message)
result = call_tool(tools, name, arguments)
messages.append({
    "role": "tool",
    "tool_name": name,
    "content": json.dumps(result, ensure_ascii=False),
})
\end{Code}

\nolinkurl{assistant_message} 要保留 \nolinkurl{tool_calls} 等完整字段，不能因为 content 为空就仅记录空文本。工具结果是字典，Ollama 此处的 content 使用字符串，因此通过 json.dumps 编码。

Python 字典、JSON 文本、包含 JSON 文本的外层消息是不同层次。打印时出现转义引号和 \code{\n}，是字符串表示，并不意味着文件正文被改写。

本讲采用 Ollama 的 \nolinkurl{tool_name} 格式；不能据此推断所有供应商都使用相同字段。调用 id 随原始 assistant 消息保留，其他服务如何关联调用要以其文档为准。

\subsection{三个时刻不能混淆}

\begin{Code}
call_tool(...)             工具已经执行
messages.append(result)    结果只保存在本地列表
chat(messages, ...)        历史真正发给模型
\end{Code}

学习者已准确解释：append 不会发送网络请求，要等下一轮 chat 才把工具结果发给模型。chat 是函数对象，在本实验中的实现使用 urlopen 发起真实 HTTP 请求；测试中的 chat 则可能只是预设响应函数。

\section{从 mock 迁移到真实自动循环}

\subsection{两种协议的差异}

{\small\setlength{\tabcolsep}{4pt}
\renewcommand{\arraystretch}{1.25}
\begin{longtable}{>{\raggedright\arraybackslash}p{0.23\linewidth}>{\raggedright\arraybackslash}p{0.33\linewidth}>{\raggedright\arraybackslash}p{0.33\linewidth}}
\toprule
\textbf{项目} & \textbf{l06 mock} & \textbf{原生 Ollama 循环} \\
\midrule\endfirsthead
\toprule
\textbf{项目} & \textbf{l06 mock} & \textbf{原生 Ollama 循环} \\
\midrule\endhead
\bottomrule\endfoot
模型来源 & 手写确定性 Python 分支 & 本地真实模型 HTTP 响应 \\
请求形状 & type=\nolinkurl{tool_call}、name、arguments & message.\nolinkurl{tool_calls} \\
结束标记 & type=final & 无新工具请求，有完整非空回答 \\
工具观察 & 内部教学字典 & role=tool，content 为 JSON 字符串 \\
\end{longtable}}

原有 loop.py 用来保留内部协议教学示范；\nolinkurl{ollama_loop}.py 使用原生协议。mock 的“看到错误后修正”是预写逻辑，不能当成模型自主纠错证据。

\subsection{核心循环的控制流}

\begin{Code}
messages = initial_messages
for step in range(1, max_steps + 1):
    response = chat(messages, tool_descriptions)
    message = response["message"]
    messages.append(message)
    calls = message.get("tool_calls") or []
    if calls:
        # 校验、执行所有请求，保存观察，应用策略
        continue
    # 没有调用且有完整非空文本，返回 completed
\end{Code}

这是简化教学代码；实际实现还检查响应类型、done、\nolinkurl{done_reason}、调用结构与网络错误，并保存诊断信息。

\begin{itemize}
\item messages 在循环外创建，历史才能跨轮保留。
\item continue 结束当前一轮，进入下一轮请求，不清空历史。
\item return 结束整个循环函数。
\end{itemize}

\subsection{为什么必须先看 tool\_calls}

l14 组合任务第二轮实际同时包含“现在我将读取该文件的内容”和 \nolinkurl{read_text_file} 请求。程序先执行工具，再把结果发给下一轮，不能把这段说明性文字当作最终回答。

如果删掉执行工具之后的 continue：工具可能已经执行，但随后仍会检查原 assistant 消息的 content；非空时提前返回 completed，模型没收到新观察。若 content 为空，则走到 \nolinkurl{empty_model_response}，也会提前结束。

\subsection{max\_steps 计数什么}

一步是一次模型请求。一次响应里执行两个工具，仍只消耗一次模型决策；最终回答也要消耗一次决策。

{\small\setlength{\tabcolsep}{4pt}
\renewcommand{\arraystretch}{1.25}
\begin{longtable}{>{\raggedright\arraybackslash}p{0.28\linewidth}>{\raggedright\arraybackslash}p{0.63\linewidth}}
\toprule
\textbf{轨迹} & \textbf{模型决策数} \\
\midrule\endfirsthead
\toprule
\textbf{轨迹} & \textbf{模型决策数} \\
\midrule\endhead
\bottomrule\endfoot
提出读取请求 → 根据结果回答 & 2 \\
错参数 → 修正参数 → 最终回答 & 3 \\
查看目录 → 读取发现的文件 → 回答 & 3 \\
\end{longtable}}

学习者曾将 mock 的 \nolinkurl{max_steps} 改为 2：第二次工具读取成功，但没有第三次 final 决策机会，程序正确返回 \nolinkurl{max_steps_exceeded}。\nolinkurl{max_steps} 是上限，不要求用完。

\section{失败后，修正、重试还是停止}

\subsection{先读证据，区分层次}

最初的实验诊断题中，status、\nolinkurl{exit_code}、stderr 分别表示工具总结状态、子进程退出信息和诊断输出。未来执行实验时应共同记录，还要保存实际配置、cwd、种子、数据划分和代码版本。

目前读取工具没有启动训练子进程。工具返回 status=failed 后，演示脚本可能仍以退出码 0 正常结束，因为错误已被处理。不能把工具状态和外层进程退出码当作同一个字段。

\subsection{本课的失败策略}

{\small\setlength{\tabcolsep}{4pt}
\renewcommand{\arraystretch}{1.25}
\begin{longtable}{>{\raggedright\arraybackslash}p{0.23\linewidth}>{\raggedright\arraybackslash}p{0.33\linewidth}>{\raggedright\arraybackslash}p{0.33\linewidth}}
\toprule
\textbf{工具结果} & \textbf{动作} & \textbf{理由或限制} \\
\midrule\endfirsthead
\toprule
\textbf{工具结果} & \textbf{动作} & \textbf{理由或限制} \\
\midrule\endhead
\bottomrule\endfoot
success & continue & 继续把观察交给模型 \\
\nolinkurl{invalid_arguments}，尚有预算 & correct & 允许模型根据真实错误重提参数 \\
\nolinkurl{invalid_arguments}，预算用完 & stop & \nolinkurl{correction_limit_exceeded} \\
\nolinkurl{file_not_found} & stop & 当前指定输入无法读取 \\
\nolinkurl{path_outside_project} & stop & 超出项目范围 \\
\nolinkurl{unknown_tool} & stop & 程序未注册该能力 \\
未分类失败 & stop & 尚无证据支持继续尝试 \\
\end{longtable}}

这些规则针对当前只读任务，不是所有 Agent 的通用策略。未来若允许目录搜索，缺文件可能触发受控搜索；未知工具名也可设计成一次受限修正。当前实现没有这些恢复分支。

\subsection{参数修正与原样重试}

\begin{Code}
参数修正：path → invalid_arguments → 模型改为 relative_path
原样重试：相同请求因暂时性故障失败，等待后重新提交
\end{Code}

前者改变请求内容，后者通常保持请求内容。是否适合后者，要确认错误是否可能恢复、重复执行是否安全，以及次数和时间预算是否足够。

本项目尚未实现网络自动退避、抖动或有副作用操作的幂等控制。模型请求超时仍停止，不因出现 timeout 就盲目重新提交。

\subsection{两个预算同时生效}

\nolinkurl{max_steps} 限制任务总模型请求次数；\nolinkurl{max_corrections} 默认只给一次参数修正机会，按整个任务累计，成功后不补满。

即使 \nolinkurl{max_steps}=5 还有剩余，只要第二次参数仍错、修正机会已耗尽，就停止。学习者已经独立解释这个判断。允许一次机会不保证模型修正成功，也不保证总步数一定够完成最终回答。

\section{Python 异常如何进入策略}

\subsection{抛异常与返回错误字典}

函数正常返回 failed 字典时，调用者拿得到 result；函数抛异常而未捕获时，result 赋值没有完成，后面的消息追加和策略判断都还没有执行。

分发器只在实际调用位置捕获已知异常：

\begin{Code}
try:
    return tool(**arguments)
except UnicodeDecodeError as error:
    return _exception_result(name, "invalid_encoding", error)
\end{Code}

except ... as error 将异常对象绑定到 error。type(error).\nolinkurl{__name__} 获取原始类型名；str(error) 获取说明文本。辅助函数返回失败字典后，循环才可以继续处理这个观察。

\subsection{三种职责}

{\small\setlength{\tabcolsep}{4pt}
\renewcommand{\arraystretch}{1.25}
\begin{longtable}{>{\raggedright\arraybackslash}p{0.28\linewidth}>{\raggedright\arraybackslash}p{0.63\linewidth}}
\toprule
\textbf{位置} & \textbf{负责什么} \\
\midrule\endfirsthead
\toprule
\textbf{位置} & \textbf{负责什么} \\
\midrule\endhead
\bottomrule\endfoot
except / 异常转换辅助函数 & 捕获已知异常，描述发生了什么 \\
\nolinkurl{decide_tool_result} & 根据类别和预算判断下一步动作 \\
\nolinkurl{run_ollama_loop} & 保存结果，并执行 continue 或 return \\
\end{longtable}}

策略返回 stop 只是一个动作决定，循环中的 return 才真正结束函数。except 本身不选择“哪一个策略函数”；本实验启用哪个策略，是调用循环时显式传入的。

\subsection{字段与映射}

{\small\setlength{\tabcolsep}{4pt}
\renewcommand{\arraystretch}{1.25}
\begin{longtable}{>{\raggedright\arraybackslash}p{0.23\linewidth}>{\raggedright\arraybackslash}p{0.33\linewidth}>{\raggedright\arraybackslash}p{0.33\linewidth}}
\toprule
\textbf{Python 异常} & \textbf{\nolinkurl{error_type}} & \textbf{当前策略 reason} \\
\midrule\endfirsthead
\toprule
\textbf{Python 异常} & \textbf{\nolinkurl{error_type}} & \textbf{当前策略 reason} \\
\midrule\endhead
\bottomrule\endfoot
FileNotFoundError & \nolinkurl{file_not_found} & \nolinkurl{required_file_missing} \\
PermissionError & \nolinkurl{permission_denied} & \nolinkurl{file_access_denied} \\
IsADirectoryError & \nolinkurl{is_a_directory} & \nolinkurl{expected_file} \\
NotADirectoryError & \nolinkurl{not_a_directory} & \nolinkurl{invalid_parent_directory} \\
UnicodeDecodeError & \nolinkurl{invalid_encoding} & \nolinkurl{utf8_required} \\
TimeoutError & \nolinkurl{tool_timeout} & \nolinkurl{tool_timeout_needs_review} \\
其他 OSError & \nolinkurl{io_error} & \nolinkurl{io_failure} \\
\end{longtable}}

\nolinkurl{error_type} 是程序定义的类别，供策略分支使用；\nolinkurl{exception_type} 是实际 Python 类型，方便排错；message 保存细节，策略不靠文本匹配来决定动作。

已有 \nolinkurl{read_text_file} 内部处理 FileNotFoundError 并直接返回原错误字典，因此不一定包含 \nolinkurl{exception_type}。只有传播到分发器、在这里被捕获的异常，才增加这类诊断字段。

\subsection{具体异常应先于父类捕获}

\begin{Code}
OSError
  + FileNotFoundError
  + PermissionError
  + IsADirectoryError
  + NotADirectoryError
  + TimeoutError

ValueError
  + UnicodeError
      + UnicodeDecodeError
\end{Code}

TimeoutError 是 OSError 的子类，匹配到第一个适用的 except 后就执行该分支。如果 OSError 在前，超时会被统一归为 \nolinkurl{io_error}。学习者已独立增加超时分支、策略理由，并正确解释这个顺序。

没有使用 except Exception 把所有问题都压成工具失败。工具内部意外的 TypeError 等编程错误仍会暴露，不能误报成模型参数错误；KeyboardInterrupt、SystemExit 也不在本次捕获范围。

\subsection{实际异常与注入异常要分清}

l13 用真实非法 UTF-8 字节触发 UnicodeDecodeError；目录读取在本机 Windows 上实际触发 PermissionError；另一个权限案例是工具主动 raise PermissionError，不改 Windows ACL。

目录返回 PermissionError 不足以断言发生了 ACL 配置错误，不同操作系统可能给出不同异常。临时夹具会清理，报告中的临时路径在实验后不存在是正常现象。

\section{第二个工具：目录查看与真实选择}

\subsection{目录观察的数据结构}

\begin{Code}
{
  "status": "success",
  "entries": [{
    "name": "hello.txt",
    "type": "file",
    "relative_path": "labs/fixtures/hello.txt"
  }],
  "truncated": false
}
\end{Code}

这只说明名称、类型和项目相对位置，没有正文。学习者最新回答正确：下一步把条目中的 \nolinkurl{relative_path} 放进 \nolinkurl{read_text_file} 的参数。

\subsection{目录工具的边界}

\begin{itemize}
\item 使用 iterdir 获取直接子项，不递归遍历整棵目录树。
\item 先 resolve 并检查项目范围，再枚举。
\item 按名称排序，最多返回 50 项，超过则 truncated=true。
\item 空目录是成功的空结果，不是异常。
\item 返回项目相对路径，方便接给读取工具。
\item 对符号链接标记类型，后续访问仍需检查实际解析位置。
\end{itemize}

50 项仅限制返回内容，不是枚举耗时或内存的硬上限。truncated=true 时，不能断言已经知道目录里全部文件。

\subsection{工具存在于两个层面}

Python 注册表让程序能够执行函数；请求中的工具说明让模型知道工具名称、用途和输入规则。只增加其中一个层面，不等于完成接入。

\begin{Code}
{
    "read_text_file": read_project_file,
    "list_directory": list_project_directory,
}
\end{Code}

两个工具共享 \nolinkurl{relative_path} 的类型、必填与额外键校验，函数名已通用化为 \nolinkurl{validate_relative_path_arguments}。运行时仍用模型返回的 function.name 分发，没有按任务关键词写路由规则。

\subsection{三种真实任务与预期轨迹}

{\small\setlength{\tabcolsep}{4pt}
\renewcommand{\arraystretch}{1.25}
\begin{longtable}{>{\raggedright\arraybackslash}p{0.23\linewidth}>{\raggedright\arraybackslash}p{0.33\linewidth}>{\raggedright\arraybackslash}p{0.33\linewidth}}
\toprule
\textbf{任务} & \textbf{实际工具序列} & \textbf{模型决策数} \\
\midrule\endfirsthead
\toprule
\textbf{任务} & \textbf{实际工具序列} & \textbf{模型决策数} \\
\midrule\endhead
\bottomrule\endfoot
查看目录名称和类型 & \nolinkurl{list_directory} & 2 \\
读取已知文件正文 & \nolinkurl{read_text_file} & 2 \\
先找目录中的 .txt 文件，再读取 & \nolinkurl{list_directory} → \nolinkurl{read_text_file} & 3 \\
\end{longtable}}

每个任务是新会话，向模型提供相同的两份工具说明。EXPECTED 只在运行后对照这几个任务的简单序列，不影响模型选择。序列相符也不替代对回答内容的检查。

\section{本阶段实验与结果总账}

\subsection{课程、小实验与文件对应}

{\small\setlength{\tabcolsep}{4pt}
\renewcommand{\arraystretch}{1.25}
\begin{longtable}{>{\raggedright\arraybackslash}p{0.18\linewidth}>{\raggedright\arraybackslash}p{0.29\linewidth}>{\raggedright\arraybackslash}p{0.16\linewidth}>{\raggedright\arraybackslash}p{0.22\linewidth}}
\toprule
\textbf{小实验} & \textbf{主要知识} & \textbf{对应小课} & \textbf{类型} \\
\midrule\endfirsthead
\toprule
\textbf{小实验} & \textbf{主要知识} & \textbf{对应小课} & \textbf{类型} \\
\midrule\endhead
\bottomrule\endfoot
\nolinkurl{l01_read_file} & 文件读取、UTF-8、成功与缺文件结果 & 02 & 真实文件工具 \\
\nolinkurl{l02_path_scope} & resolve、路径范围、存在与许可 & 03 & 路径观察 \\
\nolinkurl{l03_tool_registry} & 名称分发、函数引用、闭包根目录 & 04 & 手工请求与真实工具 \\
\nolinkurl{l04_tool_arguments} & 缺字段、错类型、空白、额外参数 & 05 & 参数错误实验 \\
\nolinkurl{l05_tool_schema} & 向模型描述输入契约 & 06 & 无模型 API \\
\nolinkurl{l06_mock_agent_loop} & 观察反馈、修正、最终回答、步数 & 07 & 确定性替身与真实工具 \\
\nolinkurl{l07_local_model_info} & 服务与模型元数据 & 08 & 只查询，不推理 \\
\nolinkurl{l08_first_local_chat} & 第一条真实聊天请求 & 09 & 真实模型 \\
\nolinkurl{l09_first_tool_call} & 原生 \nolinkurl{tool_calls}，仅提出请求 & 10 & 真实模型，不执行工具 \\
\nolinkurl{l10_tool_result_reply} & 执行保存请求、回传真实结果 & 11 & 真实模型与工具 \\
\nolinkurl{l11_real_agent_loop} & 自动请求、执行、回传、停止 & 12 & 真实自动循环 \\
\nolinkurl{l12_tool_failure_policy} & 错误分类与一次修正机会 & 13 & 预设请求与真实工具 \\
\nolinkurl{l13_tool_exceptions} & 原始异常、字典转换、策略停止 & 14 & 真实文件与明确注入 \\
超时分类独立练习 & 子类捕获顺序、策略映射 & 15 & 学习者独立修改 \\
\nolinkurl{l14_multi_tool_agent} & 两种工具选择与观察串联 & 16 & 三个真实模型任务 \\
\end{longtable}}

第 01 小课的 train.csv 是纸面诊断；未因此宣称训练已执行。小课编号与 16 周路线的周次不是同一个维度。

\subsection{学习者已提供的真实模型结果}

{\small\setlength{\tabcolsep}{4pt}
\renewcommand{\arraystretch}{1.25}
\begin{longtable}{>{\raggedright\arraybackslash}p{0.18\linewidth}>{\raggedright\arraybackslash}p{0.29\linewidth}>{\raggedright\arraybackslash}p{0.16\linewidth}>{\raggedright\arraybackslash}p{0.22\linewidth}}
\toprule
\textbf{实验} & \textbf{决策 / 工具观察} & \textbf{耗时} & \textbf{证据性质} \\
\midrule\endfirsthead
\toprule
\textbf{实验} & \textbf{决策 / 工具观察} & \textbf{耗时} & \textbf{证据性质} \\
\midrule\endhead
\bottomrule\endfoot
l08 & 单轮聊天成功 & 51.230 秒 & 学习者独立成功重试 \\
l09 & 1 条正确工具请求，执行数为 0 & 74.809 秒 & 学习者独立运行 \\
l10 & 读取 1 次，无新工具请求，正确回答 & 72.302 秒 & 学习者独立运行 \\
l11 & 决策 2 次，读取 1 次，正确回答 & 0.858 秒 & 学习者独立运行 \\
l14 list & \nolinkurl{list_directory}，决策 2 次 & 51.419 秒 & 学习者最新独立运行 \\
l14 read & \nolinkurl{read_text_file}，决策 2 次 & 0.856 秒 & 学习者最新独立运行 \\
l14 discover & 两工具依次调用，决策 3 次 & 1.413 秒 & 学习者最新独立运行 \\
\end{longtable}}

上表来自终端证据和对应报告，最后三项已核对本次附件。报告文件可能被同名的新运行覆盖，因此讲义会保存当前快照、原始终端附件与文件哈希。

\subsection{冷加载与热运行不能直接比较}

前一轮 AI 对 l14 三项实测耗时为 144.908、193.107、294.652 秒，多轮加载约 75–79 秒且没有输入缓存命中。学习者最新 list 的首轮加载约 31.748 秒，随后加载约 0.003 秒；read 和 discover 的模型已加载，多轮还命中输入缓存。

这解释了观测条件明显不同，但不是受控性能实验。不能据此声称代码优化了几百倍，也不能把增大 socket timeout 说成加速模型。timeout 改变允许等待的时长，不直接提高计算速度。

l08 初版曾 60 秒超时；一次 180 秒诊断请求成功约 45.028 秒。磁盘后来出现早期 60 秒版，因此要区分实际使用的版本。后续 l09–l14 的新入口独立使用 180 秒每请求等待；这是 socket timeout，不是整个任务的总截止时间。

\section{测试、回归与验收范围}

\subsection{测试数量如何增长}

{\small\setlength{\tabcolsep}{4pt}
\renewcommand{\arraystretch}{1.25}
\begin{longtable}{>{\raggedright\arraybackslash}p{0.23\linewidth}>{\raggedright\arraybackslash}p{0.33\linewidth}>{\raggedright\arraybackslash}p{0.33\linewidth}}
\toprule
\textbf{阶段} & \textbf{tests 中的方法数} & \textbf{增加的主要行为} \\
\midrule\endfirsthead
\toprule
\textbf{阶段} & \textbf{tests 中的方法数} & \textbf{增加的主要行为} \\
\midrule\endhead
\bottomrule\endfoot
文件读取 & 2 & 已有 UTF-8 文件、项目内缺失文件 \\
路径范围 & 5 & 相对越界、绝对越界、范围检查顺序 \\
注册表 & 9 & 名称、根目录绑定、分发后的路径边界 \\
参数契约 & 15 & 类型、必填、额外键、空白、文件名空格 \\
mock 循环 & 20 & 观察反馈、终止、非法响应、步数上限 \\
原生循环 & 29 & 原生消息、批次、网络与截断响应 \\
失败策略 & 35 & 修正预算、立即停止、跳过剩余工具 \\
异常转换 & 41 & 原始异常、分类、策略停止、编程错误边界 \\
目录工具 & 50 & 一层目录、参数、范围、截断与观察串联 \\
\end{longtable}}

另有 exercises 中 4 项超时分类验收。当前 AI 已验收项目 50 项与练习 4 项通过；学习者此前已独立提供 41 项通过输出，超时练习修改由 AI 验收。本次附件没有新增 50 项测试终端输出，不把它自动记为学习者独立复跑过。

unittest 的数量指测试方法，部分方法包含多个 subTest。50 项通过只支持这些覆盖行为，不等于覆盖所有故障，也不是 50 个真实 Agent 任务成功。

\subsection{两组命令都需要知道}

\begin{Code}
python -X utf8 -m unittest discover -s tests -v
python -X utf8 -m unittest discover -s exercises -p test_timeout_mapping.py -v
\end{Code}

只运行 tests 不会执行 exercises 中的独立练习。测试中可以用预设响应稳定制造错误，但真正需要验证的文件读取、参数校验和循环控制保持真实代码。

\subsection{对文件版本回退保持证据意识}

学习过程中曾出现路径检查缺失、超时分支顺序变化、目录注册项回到旧版本等现象。只读检查和回归测试发现了问题，并恢复已验收的行为；来源没有确定，不能猜测是谁覆盖了文件。

复盘时应核对 sys.executable、模块 \nolinkurl{__file__}、实际源码、文件版本以及失败测试。当前讲义附带源码快照与 SHA-256 清单，便于以后区分“当时的实现”和“后续改过的实现”。

\section{易错点速查}

{\small\setlength{\tabcolsep}{4pt}
\renewcommand{\arraystretch}{1.25}
\begin{longtable}{>{\raggedright\arraybackslash}p{0.28\linewidth}>{\raggedright\arraybackslash}p{0.63\linewidth}}
\toprule
\textbf{容易混淆的说法} & \textbf{正确理解} \\
\midrule\endfirsthead
\toprule
\textbf{容易混淆的说法} & \textbf{正确理解} \\
\midrule\endhead
\bottomrule\endfoot
注册表里只有名称字符串 & 键是名称字符串，值是函数对象 \\
写函数名就是调用 & 加括号并传参才执行；不加括号可传递函数对象 \\
模型生成 \nolinkurl{tool_calls} 就已读文件 & 只提出请求，Python 还要校验并执行 \\
content 为空就是失败 & 可能正在请求工具，要同时检查 \nolinkurl{tool_calls} \\
done=true 表示任务结束 & 只表示本轮生成完成 \\
append 结果就已发给模型 & 只是本地保存，要等下一次 chat \\
有一段文字就可结束 & 有 \nolinkurl{tool_calls} 时先处理工具请求 \\
\nolinkurl{max_steps} 有剩余就继续 & 修正预算或明确停止规则可提前结束 \\
except 负责决定停止策略 & except 转换；策略判断；循环执行终止 \\
捕获 TypeError 就算参数校验 & 校验应发生在执行前，工具内部 bug 不能误归类 \\
目录名字已经代表文件内容 & 名称、类型、路径是元信息，正文还需读取 \\
minLength=1 能排除全空格 & 空格有长度，需要非空白约束 \\
strip 删除所有空格 & 只处理两端空白；本程序仅用它判断 \\
文件范围检查等于完整沙箱 & 只实现当前路径约束，隔离与资源限制仍待学 \\
一次热运行更快证明优化 & 条件不同，需区分加载、缓存和生成时间 \\
测试通过表示整个 Agent 已可靠 & 只支持已覆盖行为，还需独立任务评测 \\
\end{longtable}}

\section{自测题与参考答案}

\subsection{先不看答案，尝试复述}

\begin{enumerate}
\item 工具说明、注册表、工具函数分别解决什么问题？
\item 给定 tools 字典，取出函数与执行函数分别是哪一行？
\item 项目外文件不存在，先返回什么错误？为什么？
\item \code{   a   } 是否通过当前字符串校验？它是否会自动变成 a？
\item 模型有 \nolinkurl{tool_calls} 且 content 非空时，应先处理什么？
\item 文件工具成功后，为什么还需要一次模型请求？
\item append、continue、chat、return 分别做什么？
\item 第二次参数仍错，但 \nolinkurl{max_steps} 还有剩余，如何决定？
\item \nolinkurl{error_type}、\nolinkurl{exception_type} 和 \nolinkurl{stop_reason} 分别代表什么？
\item TimeoutError 放到 OSError 后有什么后果？
\item 目录列出 hello.txt 后，下一次读取使用哪个字段？
\item 为什么不能用 0.856 秒与 193.107 秒直接宣称优化效果？
\end{enumerate}

\subsection{参考答案}

\begin{enumerate}
\item 工具说明告诉模型能做什么及参数契约；注册表把名称映射为函数；函数实际接触环境。
\item \code{tool = tools[name]} 取出对象；\code{tool(**arguments)} 执行。
\item \nolinkurl{path_outside_project}。范围检查先于环境访问，不用文件是否存在来决定许可。
\item 可以通过；不会自动变成 a，本程序不改写原路径。
\item 先执行请求并保存观察，不能将说明性文字当作最终答案。
\item 工具结果尚需发送给模型，模型才能根据观察生成回答或下一步动作。
\item append 本地保存；continue 进入下一轮；本实验 chat 发网络请求；return 结束函数。
\item 修正机会已用完则停止，两个预算同时约束任务。
\item \nolinkurl{error_type} 是程序类别，\nolinkurl{exception_type} 是原始 Python 类型，\nolinkurl{stop_reason} 是策略停止理由。
\item 被父类分支先接住，可能失去独立超时类别。
\item 条目里的 \nolinkurl{relative_path}；名称、类型和路径都不包含正文。
\item 模型加载、缓存命中等条件不同，未控制变量，样本也不足。
\end{enumerate}

\subsection{已完成的独立修改}

学习者在 registry.py 中增加了 TimeoutError 的专门分支，在 \nolinkurl{tool_policy}.py 中增加 \nolinkurl{tool_timeout_needs_review}；能够解释父子异常的捕获顺序。4 项练习与原有回归通过。

这是已完成的一次独立功能扩展。尚未要求从空白编写完整注册表、原生循环或第二个工具；下次可从读代码转向小范围独立实现。

\section{阶段结论与下次恢复入口}

\subsection{当前进度}

当前已经有最小 Agent 的工程原型，并具备几种已验证的失败终态。学习者能够解释核心反馈关系，并完成一次实际功能修改。第 1–3 周仍有未完成的任务清单、课程阅读和独立实现验收，不按小课数量直接推算总进度百分比。

{\small\setlength{\tabcolsep}{4pt}
\renewcommand{\arraystretch}{1.25}
\begin{longtable}{>{\raggedright\arraybackslash}p{0.28\linewidth}>{\raggedright\arraybackslash}p{0.63\linewidth}}
\toprule
\textbf{范围} & \textbf{当前状态} \\
\midrule\endfirsthead
\toprule
\textbf{范围} & \textbf{当前状态} \\
\midrule\endhead
\bottomrule\endfoot
工具读取、注册、参数与 Schema & 已学习并独立运行，多处机制已复述 \\
本地模型与原生工具循环 & 已独立跑通 \\
修正预算与异常处理 & 已运行，超时分支已独立修改 \\
双工具选择与目录观察接读取 & 最新三项独立运行成功，字段含义回答正确 \\
整个循环与工具的独立实现 & 待验收 \\
CS2680 讲义、论文逐项阅读 & 尚未逐项验收 \\
检索、记忆、规划、独立评测 & 后续阶段 \\
受控实验执行与真实指标验证 & 毕业项目下一重要阶段 \\
\end{longtable}}

\subsection{下一步路线，仅作恢复提示}

\begin{enumerate}
\item 不看答案复述“任务→请求→校验→工具→观察→下一轮”，并看懂两工具的消息记录。
\item 做一次由学习者亲自编写的小功能或参数规则扩展。
\item 增加项目代码定位与受控实验执行，再读取本次运行的指标。
\item 首个实验目标仍为小型合成数据线性回归：固定条件比较两个学习率，报告验证集 MSE 和耗时。该训练实验尚未创建或执行。
\item 后续依次补齐检索、记忆恢复、规划对照、独立评测、权限隔离、框架/MCP、成本优化与模型服务实验。
\end{enumerate}

今天到此收束，不在本次归档中新增 Agent 功能。下次可以从本讲自测题和 PROGRESS.md 接续，不必重新安装环境或重复全部旧实验。

\section{参考入口与证据使用说明}

以下为本项目已保存或此前访问过的资料入口，不代表学习者已完整阅读。接口说明以项目实际采用的 Ollama 格式为例，其他服务需要适配。

\begin{itemize}
\item \href{https://cs2680.com/}{CS2680 正式课程入口}
\item \href{https://leaderboard.cs2680.com/}{CS2680 排行榜入口}
\item \href{https://huggingface.co/learn/agents-course/en/unit0/introduction}{Hugging Face Agents Course}
\item \href{https://llmagents-learning.org/f24}{Berkeley LLM Agents F24}
\item \href{https://llmagents-learning.org/sp25}{Berkeley Advanced LLM Agents SP25}
\item \href{https://docs.ollama.com/api/chat}{Ollama Chat API}
\item \href{https://docs.ollama.com/capabilities/tool-calling}{Ollama Tool Calling}
\item \href{https://docs.python.org/3.11/library/pathlib.html}{Python 3.11 pathlib}
\item \href{https://docs.python.org/3.11/library/exceptions.html}{Python 3.11 内置异常}
\item \href{https://json-schema.org/understanding-json-schema/reference/object}{JSON Schema 对象规则}
\item \href{https://json-schema.org/understanding-json-schema/reference/string}{JSON Schema 字符串规则}
\item \href{https://aws.amazon.com/builders-library/timeouts-retries-and-backoff-with-jitter/}{AWS 超时、重试与退避}
\end{itemize}

正文解释来自本项目实际代码、16 次小课、运行笔记、原始 JSON 报告和学习者终端证据。后附完整文件索引、关键源码快照和运行报告清单，便于复核。

\section{附录 A. 全部学习资料索引与快照}
本次固定保存 95 个原始文件。下列链接指向讲义目录内 \nolinkurl{source_snapshot}，不随原项目后续编辑而变化。manifest.json 记录字节数和 SHA-256；它保证归档内容可核对，不证明课程已全部掌握。

\subsection{项目入口与环境}
\begin{itemize}
\item \nolinkurl{.gitignore}
\item \nolinkurl{PROGRESS.md}：学习进度与验收
\item \nolinkurl{README.md}：Agent 学习与构建项目
\item \nolinkurl{environment.yml}
\end{itemize}

\subsection{学习规划与资料}
\begin{itemize}
\item \nolinkurl{docs/01-roadmap.md}：16 周学习路线
\item \nolinkurl{docs/02-knowledge-map.md}：Agent 知识点覆盖矩阵
\item \nolinkurl{docs/03-resources.md}：课程与资料导航
\item \nolinkurl{docs/04-labs.md}：递进实验任务书
\item \nolinkurl{docs/05-capstone.md}：毕业项目：ExperimentAgent 代码实验助手
\item \nolinkurl{docs/06-learning-workflow.md}：学习方法与环境
\end{itemize}

\subsection{逐次讲义}
\begin{itemize}
\item \nolinkurl{lessons/01-agent-loop.md}：第 1 课：代码实验 Agent 如何工作
\item \nolinkurl{lessons/02-first-file-tool.md}：第 2 次动手：文件读取工具与 Anaconda
\item \nolinkurl{lessons/03-path-scope.md}：第 3 次动手：工具可以读哪些路径
\item \nolinkurl{lessons/04-tool-registry.md}：第 4 次动手：根据工具名称调用函数
\item \nolinkurl{lessons/05-tool-arguments.md}：第 5 次动手：调用工具前检查参数
\item \nolinkurl{lessons/06-tool-json-schema.md}：第 6 次动手：用 JSON Schema 描述工具参数
\item \nolinkurl{lessons/07-mock-agent-loop.md}：第 7 次动手：把工具接进反馈循环
\item \nolinkurl{lessons/08-model-backend.md}：第 8 次动手：选择真实模型接口
\item \nolinkurl{lessons/09-first-local-chat.md}：第 9 次动手：第一次真实模型请求
\item \nolinkurl{lessons/10-real-tool-call.md}：第 10 次动手：观察真实模型的工具调用请求
\item \nolinkurl{lessons/11-tool-result-reply.md}：第 11 次动手：执行工具并回传结果
\item \nolinkurl{lessons/12-real-agent-loop.md}：第 12 次动手：真实模型的自动循环
\item \nolinkurl{lessons/13-tool-failure-policy.md}：第 13 次动手：工具失败后，修正还是停止
\item \nolinkurl{lessons/14-tool-exceptions.md}：第 14 次动手：把 Python 异常转换成工具结果
\item \nolinkurl{lessons/15-timeout-mapping-exercise.md}：第 15 次动手：由你完成工具超时分类
\item \nolinkurl{lessons/16-directory-tool-and-selection.md}：第 16 次动手：第二个工具与模型选择
\item \nolinkurl{lessons/week01.md}：第 1 周：理解代码实验 Agent，并定义任务
\end{itemize}

\subsection{运行与反馈记录}
\begin{itemize}
\item \nolinkurl{notes/2026-10-21-review.md}：2026-10-21 阶段复盘
\item \nolinkurl{notes/argument-validation-run.md}：参数校验实操与反馈
\item \nolinkurl{notes/directory-tools-run.md}：查看目录与双工具选择实测
\item \nolinkurl{notes/file-tool-run.md}：文件读取工具实测记录
\item \nolinkurl{notes/first-local-chat-run.md}：第一次本地模型聊天实测
\item \nolinkurl{notes/first-native-tool-call.md}：真实模型首次原生工具请求
\item \nolinkurl{notes/lesson01-feedback.md}：第 1 课：首次练习记录
\item \nolinkurl{notes/local-chat-timeout.md}：本地聊天超时诊断与修正
\item \nolinkurl{notes/mock-loop-run.md}：确定性模型替身闭环首次运行
\item \nolinkurl{notes/path-scope-run.md}：路径范围检查：运行与反馈
\item \nolinkurl{notes/real-agent-loop-run.md}：真实模型自动循环实测
\item \nolinkurl{notes/timeout-mapping-exercise.md}：独立修改练习准备记录
\item \nolinkurl{notes/tool-exceptions-run.md}：工具异常转换实测
\item \nolinkurl{notes/tool-failure-policy-run.md}：工具失败策略实验记录
\item \nolinkurl{notes/tool-registry-run.md}：工具注册表首次示范与运行
\item \nolinkurl{notes/tool-result-reply-run.md}：执行真实工具请求并回传结果
\item \nolinkurl{notes/tool-schema-run.md}：工具参数 JSON Schema 首次演示
\end{itemize}

\subsection{运行实验与夹具}
\begin{itemize}
\item \nolinkurl{labs/__init__.py}
\item \nolinkurl{labs/fixtures/hello.txt}
\item \nolinkurl{labs/l01_read_file.py}
\item \nolinkurl{labs/l02_path_scope.py}
\item \nolinkurl{labs/l03_tool_registry.py}
\item \nolinkurl{labs/l04_tool_arguments.py}
\item \nolinkurl{labs/l05_tool_schema.py}
\item \nolinkurl{labs/l06_mock_agent_loop.py}
\item \nolinkurl{labs/l07_local_model_info.py}
\item \nolinkurl{labs/l08_first_local_chat.py}
\item \nolinkurl{labs/l09_first_tool_call.py}
\item \nolinkurl{labs/l10_tool_result_reply.py}
\item \nolinkurl{labs/l11_real_agent_loop.py}
\item \nolinkurl{labs/l12_tool_failure_policy.py}
\item \nolinkurl{labs/l13_tool_exceptions.py}
\item \nolinkurl{labs/l14_multi_tool_agent.py}
\end{itemize}

\subsection{完整 JSON 运行报告}
\begin{itemize}
\item \nolinkurl{reports/l08_first_local_chat.json}
\item \nolinkurl{reports/l08_initial_success.json}
\item \nolinkurl{reports/l08_timeout_diagnostic.json}
\item \nolinkurl{reports/l09_first_tool_call.json}
\item \nolinkurl{reports/l10_tool_result_reply.json}
\item \nolinkurl{reports/l11_real_agent_loop_steps1.json}
\item \nolinkurl{reports/l11_real_agent_loop_steps5.json}
\item \nolinkurl{reports/l12_tool_failure_policy.json}
\item \nolinkurl{reports/l13_tool_exceptions.json}
\item \nolinkurl{reports/l14_multi_tool_agent_discover.json}
\item \nolinkurl{reports/l14_multi_tool_agent_list.json}
\item \nolinkurl{reports/l14_multi_tool_agent_read.json}
\end{itemize}

\subsection{Agent 实现与参数 Schema}
\begin{itemize}
\item \nolinkurl{experiment_agent/__init__.py}
\item \nolinkurl{experiment_agent/loop.py}
\item \nolinkurl{experiment_agent/ollama_loop.py}
\item \nolinkurl{experiment_agent/tool_policy.py}
\item \nolinkurl{experiment_agent/tools/__init__.py}
\item \nolinkurl{experiment_agent/tools/directory_tools.py}
\item \nolinkurl{experiment_agent/tools/file_tools.py}
\item \nolinkurl{experiment_agent/tools/list_directory.schema.json}
\item \nolinkurl{experiment_agent/tools/read_text_file.schema.json}
\item \nolinkurl{experiment_agent/tools/registry.py}
\end{itemize}

\subsection{项目测试}
\begin{itemize}
\item \nolinkurl{tests/test_directory_tools.py}
\item \nolinkurl{tests/test_file_tools.py}
\item \nolinkurl{tests/test_loop.py}
\item \nolinkurl{tests/test_ollama_loop.py}
\item \nolinkurl{tests/test_tool_exceptions.py}
\item \nolinkurl{tests/test_tool_failure_policy.py}
\item \nolinkurl{tests/test_tool_registry.py}
\end{itemize}

\subsection{独立练习验收}
\begin{itemize}
\item \nolinkurl{exercises/test_timeout_mapping.py}
\end{itemize}

\subsection{学习与实验记录模板}
\begin{itemize}
\item \nolinkurl{templates/experiment-report.md}：实验报告：<编号与标题>
\item \nolinkurl{templates/study-note.md}：学习记录：<主题>
\end{itemize}

\subsection{学习者原始终端附件}
\begin{itemize}
\item \nolinkurl{learner_evidence/learner_l11_output.txt}
\item \nolinkurl{learner_evidence/learner_l12_output.txt}
\item \nolinkurl{learner_evidence/learner_l14_output.txt}
\end{itemize}

\section{附录 B. 关键源码：按调用链复习}
以下摘录是本次归档时的真实源码，不是额外实现的新版本。完整实验与测试见附录 A。

\subsection{文件读取工具}

来源：\nolinkurl{experiment_agent/tools/file_tools.py}。

\begin{Code}
"""文件读取工具。"""

from pathlib import Path


def read_text_file(project_root: str, relative_path: str) -> dict:
    """仅尝试读取指定项目范围内的 UTF-8 文本。"""
    root = Path(project_root).resolve()
    path = (root / relative_path).resolve()

    if not path.is_relative_to(root):
        return {
            "status": "failed",
            "error_type": "path_outside_project",
            "path": str(path),
        }

    try:
        content = path.read_text(encoding="utf-8")
        return {
            "status": "success",
            "content": content,
        }
    except FileNotFoundError:
        return {
            "status": "failed",
            "error_type": "file_not_found",
            "path": str(path),
        }
\end{Code}

\subsection{目录工具}

来源：\nolinkurl{experiment_agent/tools/directory_tools.py}。

\begin{Code}
"""项目范围内的一层目录查看工具。"""

from pathlib import Path


def list_directory(project_root: str, relative_path: str) -> dict:
    """列出直接子项，最多返回 50 项；文件正文由读取工具获取。"""
    root = Path(project_root).resolve()
    path = (root / relative_path).resolve()
    if not path.is_relative_to(root):
        return {"status": "failed", "error_type": "path_outside_project", "path": str(path)}

    children = sorted(path.iterdir(), key=lambda child: child.name.casefold())
    entries = []
    for child in children[:50]:
        if child.is_symlink():
            entry_type = "symlink"
        elif child.is_dir():
            entry_type = "directory"
        elif child.is_file():
            entry_type = "file"
        else:
            entry_type = "other"
        entries.append({
            "name": child.name,
            "type": entry_type,
            "relative_path": child.relative_to(root).as_posix(),
        })
    return {"status": "success", "entries": entries, "truncated": len(children) > 50}
\end{Code}

\subsection{注册、参数校验与异常转换}

来源：\nolinkurl{experiment_agent/tools/registry.py}。

\begin{Code}
"""工具名称与可调用函数的映射。"""

from pathlib import Path

from experiment_agent.tools.directory_tools import list_directory
from experiment_agent.tools.file_tools import read_text_file


def build_tool_registry(project_root: str) -> dict:
    """创建绑定到指定项目根目录的工具表。"""
    root = str(Path(project_root).resolve())

    def read_project_file(relative_path: str) -> dict:
        return read_text_file(root, relative_path)

    def list_project_directory(relative_path: str) -> dict:
        return list_directory(root, relative_path)

    return {
        "read_text_file": read_project_file,
        "list_directory": list_project_directory,
    }


def validate_relative_path_arguments(arguments: object) -> str | None:
    """参数合法时返回 None，否则返回错误说明。"""
    if not isinstance(arguments, dict):
        return "arguments must be a dictionary"

    if "relative_path" not in arguments:
        return "relative_path is required"

    if set(arguments) != {"relative_path"}:
        return "only relative_path is allowed"

    value = arguments["relative_path"]
    if not isinstance(value, str) or not value.strip():
        return "relative_path must be a non-empty string"

    return None


def _exception_result(name: str, error_type: str, error: Exception) -> dict:
    """把已知执行异常转换成可记录、可交给策略判断的结果。"""
    return {
        "status": "failed",
        "error_type": error_type,
        "tool_name": name,
        "exception_type": type(error).__name__,
        "message": str(error),
    }


def call_tool(tools: dict, name: str, arguments: object) -> dict:
    """按名称查找并校验参数，在实际执行边界转换已知文件异常。"""
    tool = tools.get(name)

    if tool is None:
        return {
            "status": "failed",
            "error_type": "unknown_tool",
            "tool_name": name,
        }

    if name in {"read_text_file", "list_directory"}:
        message = validate_relative_path_arguments(arguments)
        if message is not None:
            return {
                "status": "failed",
                "error_type": "invalid_arguments",
                "message": message,
            }

    try:
        return tool(**arguments)
    except FileNotFoundError as error:
        return _exception_result(name, "file_not_found", error)
    except PermissionError as error:
        return _exception_result(name, "permission_denied", error)
    except IsADirectoryError as error:
        return _exception_result(name, "is_a_directory", error)
    except NotADirectoryError as error:
        return _exception_result(name, "not_a_directory", error)
    except UnicodeDecodeError as error:
        return _exception_result(name, "invalid_encoding", error)
    except TimeoutError as error:
        return _exception_result(name, "tool_timeout", error)
    except OSError as error:
        return _exception_result(name, "io_error", error)
\end{Code}

\subsection{策略函数}

来源：\nolinkurl{experiment_agent/tool_policy.py}。

\begin{Code}
"""本课针对只读文件任务的工具结果策略。"""


def decide_tool_result(result, corrections_used, max_corrections=1):
    """返回 continue、correct 或 stop；修正预算按整个任务累计。"""
    if not isinstance(result, dict) or result.get("status") not in {"success", "failed"}:
        return {"action": "stop", "reason": "invalid_tool_result"}
    if result["status"] == "success":
        return {"action": "continue", "reason": "tool_succeeded"}

    error_type = result.get("error_type")
    if error_type == "invalid_arguments":
        if corrections_used < max_corrections:
            return {"action": "correct", "reason": "ask_model_to_fix_arguments"}
        return {"action": "stop", "reason": "correction_limit_exceeded"}

    reasons = {
        "file_not_found": "required_file_missing",
        "path_outside_project": "permission_boundary",
        "unknown_tool": "tool_unavailable",
        "permission_denied": "file_access_denied",
        "is_a_directory": "expected_file",
        "not_a_directory": "invalid_parent_directory",
        "invalid_encoding": "utf8_required",
        "tool_timeout": "tool_timeout_needs_review",
        "io_error": "io_failure",
        
    }
    return {"action": "stop", "reason": reasons.get(error_type, "unclassified_tool_failure")}
\end{Code}

\subsection{原生消息自动循环}

来源：\nolinkurl{experiment_agent/ollama_loop.py}。

\begin{Code}
"""使用 Ollama 原生消息格式的模型—工具循环。"""

import json

from experiment_agent.tools.registry import call_tool


def run_ollama_loop(task, chat, tools, tool_descriptions, max_steps=5, tool_policy=None):
    """chat 返回原生响应；可选 tool_policy 决定工具结果出现后继续、修正或停止。"""
    if max_steps < 1:
        raise ValueError("max_steps must be at least 1")
    messages = [
        {"role": "system", "content": "你是代码实验助手。读取项目文件必须通过提供的工具获得实际内容。"},
        {"role": "user", "content": task},
    ]
    responses = []
    tool_decisions = []
    corrections_used = 0

    def finish(status, steps, **fields):
        return {
            "status": status, "steps": steps, "messages": messages,
            "responses": responses, "tool_decisions": tool_decisions, **fields,
        }

    # 一步是一次模型请求；一条响应里的多个工具调用仍属于这一步。
    for step in range(1, max_steps + 1):
        try:
            response = chat(messages, tool_descriptions)
        except (OSError, json.JSONDecodeError) as error:
            return finish("failed", step, error_type="model_request_failed", message=str(error))
        responses.append(response)

        if not isinstance(response, dict):
            return finish("failed", step, error_type="invalid_model_response")
        if response.get("done") is not True or response.get("done_reason") != "stop":
            return finish("failed", step, error_type="incomplete_model_response")
        message = response.get("message")
        if (
            not isinstance(message, dict)
            or message.get("role") != "assistant"
            or not isinstance(message.get("content", ""), str)
        ):
            return finish("failed", step, error_type="invalid_model_response")
        calls = message.get("tool_calls")
        if calls is None:
            calls = []
        if not isinstance(calls, list):
            return finish("failed", step, error_type="invalid_model_response")
        # 先检查整批调用结构，避免处理到一半才发现后一条结构不完整。
        for call in calls:
            function = call.get("function") if isinstance(call, dict) else None
            if (
                not isinstance(function, dict)
                or not isinstance(function.get("name"), str)
                or "arguments" not in function
            ):
                return finish("failed", step, error_type="invalid_model_response")

        # 保留原生 assistant 消息，包括调用 id、名称和参数。
        messages.append(message)
        if calls:
            for index, call in enumerate(calls):
                function = call["function"]
                result = call_tool(tools, function["name"], function["arguments"])
                messages.append({
                    "role": "tool",
                    "tool_name": function["name"],
                    "content": json.dumps(result, ensure_ascii=False),
                })
                if tool_policy is not None:
                    decision = tool_policy(result, corrections_used)
                    if decision["action"] == "correct":
                        corrections_used += 1
                    tool_decisions.append({
                        "step": step,
                        "tool_name": function["name"],
                        **decision,
                        "corrections_used": corrections_used,
                    })
                    if decision["action"] == "stop":
                        return finish(
                            "failed", step,
                            error_type="tool_policy_stopped",
                            tool_error_type=result.get("error_type"),
                            stop_reason=decision["reason"],
                            pending_tool_calls=len(calls) - index - 1,
                        )
            continue  # 带上真实工具结果，进入下一次模型请求。

        content = message.get("content", "")
        if content.strip():
            return finish("completed", step, answer=content)
        return finish("failed", step, error_type="empty_model_response")

    return finish("failed", max_steps, error_type="max_steps_exceeded")
\end{Code}

\subsection{内部 mock 协议循环，作为对照}

来源：\nolinkurl{experiment_agent/loop.py}。

\begin{Code}
"""通过消息历史连接模型决策和工具执行。"""

from experiment_agent.tools.registry import call_tool


def run_loop(task: str, model, tools: dict, tool_descriptions: list, max_steps: int = 5) -> dict:
    """运行内部教学协议的模型—工具循环。"""
    if max_steps < 1:
        raise ValueError("max_steps must be at least 1")

    messages = [{"role": "user", "content": task}]

    for step in range(1, max_steps + 1):
        action = model(messages, tool_descriptions)
        messages.append({"role": "assistant", "content": action})

        if isinstance(action, dict):
            if action.get("type") == "final" and isinstance(action.get("content"), str):
                return {
                    "status": "completed",
                    "answer": action["content"],
                    "steps": step,
                    "messages": messages,
                }

            if (
                action.get("type") == "tool_call"
                and isinstance(action.get("name"), str)
                and "arguments" in action
            ):
                result = call_tool(tools, action["name"], action["arguments"])
                messages.append({
                    "role": "tool",
                    "name": action["name"],
                    "content": result,
                })
                continue

        return {
            "status": "failed",
            "error_type": "invalid_model_response",
            "steps": step,
            "messages": messages,
        }

    return {
        "status": "failed",
        "error_type": "max_steps_exceeded",
        "steps": max_steps,
        "messages": messages,
    }
\end{Code}

\subsection{读取工具 Schema}

来源：\nolinkurl{experiment_agent/tools/read_text_file.schema.json}。

\begin{Code}
{
  "$schema": "https://json-schema.org/draft/2020-12/schema",
  "type": "object",
  "properties": {
    "relative_path": {
      "type": "string",
      "description": "相对于项目根目录的 UTF-8 文本文件路径，例如 labs/fixtures/hello.txt。路径须位于项目内，且不能是空字符串或全空白。",
      "minLength": 1,
      "pattern": "\\S"
    }
  },
  "required": ["relative_path"],
  "additionalProperties": false
}
\end{Code}

\subsection{目录工具 Schema}

来源：\nolinkurl{experiment_agent/tools/list_directory.schema.json}。

\begin{Code}
{
  "$schema": "https://json-schema.org/draft/2020-12/schema",
  "type": "object",
  "properties": {
    "relative_path": {
      "type": "string",
      "description": "相对于项目根目录的目录路径。使用 . 查看根目录，例如 labs/fixtures 查看该子目录。路径须位于项目内，且不能是空字符串或全空白。",
      "minLength": 1,
      "pattern": "\\S"
    }
  },
  "required": ["relative_path"],
  "additionalProperties": false
}
\end{Code}

\section{附录 C. 当前报告快照清单}
同名报告代表快照时的最近保存结果，不一定是首次运行。上文已将关键结果的学习者/AI 来源分别注明。原始 JSON 保留模型响应、消息历史或实验案例；需要精确核对时打开原文件。

{\small\setlength{\tabcolsep}{4pt}
\renewcommand{\arraystretch}{1.25}
\begin{longtable}{>{\raggedright\arraybackslash}p{0.23\linewidth}>{\raggedright\arraybackslash}p{0.33\linewidth}>{\raggedright\arraybackslash}p{0.33\linewidth}}
\toprule
\textbf{报告} & \textbf{结构与用途} & \textbf{时长或案例数} \\
\midrule\endfirsthead
\toprule
\textbf{报告} & \textbf{结构与用途} & \textbf{时长或案例数} \\
\midrule\endhead
\bottomrule\endfoot
\nolinkurl{l08_first_local_chat.json} & 模型响应 done=True & 51.230 秒 \\
\nolinkurl{l08_initial_success.json} & 模型响应 done=True & 50.340 秒 \\
\nolinkurl{l08_timeout_diagnostic.json} & 模型响应 done=True & 45.028 秒 \\
\nolinkurl{l09_first_tool_call.json} & 模型响应 done=True & 74.809 秒 \\
\nolinkurl{l10_tool_result_reply.json} & 模型响应 done=True & 72.302 秒 \\
\nolinkurl{l11_real_agent_loop_steps1.json} & 循环终态 failed & 0.522 秒 \\
\nolinkurl{l11_real_agent_loop_steps5.json} & 循环终态 completed & 0.858 秒 \\
\nolinkurl{l12_tool_failure_policy.json} & 预设请求/异常演示的逐案例结果 & 5 个案例 \\
\nolinkurl{l13_tool_exceptions.json} & 预设请求/异常演示的逐案例结果 & 4 个案例 \\
\nolinkurl{l14_multi_tool_agent_discover.json} & 循环终态 completed & 1.413 秒 \\
\nolinkurl{l14_multi_tool_agent_list.json} & 循环终态 completed & 51.419 秒 \\
\nolinkurl{l14_multi_tool_agent_read.json} & 循环终态 completed & 0.856 秒 \\
\end{longtable}}

\section{附录 D. K01–K27 知识域定位}
此处只标记本阶段覆盖程度；详细内容和未来验收标准见源文件 docs/02-knowledge-map.md。局部接触不代表整域完成。

{\small\setlength{\tabcolsep}{4pt}
\renewcommand{\arraystretch}{1.25}
\begin{longtable}{>{\raggedright\arraybackslash}p{0.23\linewidth}>{\raggedright\arraybackslash}p{0.33\linewidth}>{\raggedright\arraybackslash}p{0.33\linewidth}}
\toprule
\textbf{编号} & \textbf{知识域} & \textbf{本阶段定位} \\
\midrule\endfirsthead
\toprule
\textbf{编号} & \textbf{知识域} & \textbf{本阶段定位} \\
\midrule\endhead
\bottomrule\endfoot
K01 & Agent 基本模型 & 基础闭环已学习，独立完整实现待验收 \\
K02 & LLM 与推理 & 接触生成参数与接口，模型原理未系统复习 \\
K03 & 模型接入 & 本地聊天、原生工具消息已独立运行 \\
K04 & Prompt 与上下文 & 工具说明与消息历史局部接触 \\
K05 & 工具设计 & 工具契约已运行，超时分支独立修改 \\
K06 & 运行时控制 & 循环、停止、修正预算已实践；恢复/幂等未做 \\
K07 & 检索与知识 & 后续学习，尚未验收 \\
K08 & Agentic RAG & 后续学习，尚未验收 \\
K09 & 记忆 & 后续学习，尚未验收 \\
K10 & 规划与推理 & 后续学习，尚未验收 \\
K11 & 多 Agent & 后续学习，尚未验收 \\
K12 & 协议与框架 & 后续学习，尚未验收 \\
K13 & 协议扩展 & 后续学习，尚未验收 \\
K14 & 评测设计 & 已有单元测试与简单轨迹核对；独立任务评测未做 \\
K15 & 实验方法 & 已区分冷/热观测；尚无受控优化实验 \\
K16 & 可观测性 & 保存 JSON 报告与错误；完整追踪系统未做 \\
K17 & 安全与权限 & 已有路径边界；更完整权限和注入实验未做 \\
K18 & 执行隔离 & 后续学习，尚未验收 \\
K19 & 人机协作 & 后续学习，尚未验收 \\
K20 & 应用工程 & 专用环境与模块运行已完成；产品化未做 \\
K21 & Agent 经济性 & 接触调用次数与时间，经济性实验未做 \\
K22 & 系统基础 & 后续学习，尚未验收 \\
K23 & 模型服务 & 观察加载与缓存指标；系统服务实验未做 \\
K24 & 服务优化 & 后续学习，尚未验收 \\
K25 & 学习型 Agent & 后续学习，尚未验收 \\
K26 & 自改进系统 & 后续学习，尚未验收 \\
K27 & 特定应用 & 选定代码实验 Agent，尚未执行训练任务 \\
\end{longtable}}

\end{document}
